% GENERATED FILE. Edit canonical lessons, then run npm run generate:books.
% canonical-source-hash: fnv1a64:3758a1479e06ec55
% canonical-lessons: JA-C60-man, JA-C60-juuni, JA-C60-nasu, JA-C60-uta, JA-C60-eda

\chapter{Ten thousand, Twelve, Aubergine, Song, Branch}
\label{ch:ja-ten-thousand-twelve-aubergine-song-branch}
\hlchaptermodality{60}

\begin{chapteropening}
\textbf{By the end of this chapter:} \emph{I can say ten thousand, twelve, an aubergine, a song and a branch.}

Complete the last lesson of chapter 60: i can say ten thousand, twelve, an aubergine, a song and a branch.
\end{chapteropening}

\section[man]{\ja{まん} (man) — ten thousand}

\label{lesson:JA-C60-man}



\begin{quote}
Before the new one: say the Japanese for thirty, then the Japanese for a thousand.
\end{quote}

\subsection*{You'll want to know: \ja{まん}}
\textbf{\ja{まん}} — \emph{man} — \textquotedblleft{}ten thousand\textquotedblright{}.

\begin{grammarlens}[title={What you've built}]
One more word for this chapter.
\end{grammarlens}

\subsection*{Your turn}
\begin{itemize}
\raggedright
  \item \emph{Say it:} \emph{man}
  \item \emph{Say it:} \emph{man}, once more
  \item \emph{Say it:} say \emph{man}
  \item \emph{Recall:} say the Japanese for thirty, then the Japanese for a thousand, then say \emph{man} again
\end{itemize}

\begin{tcolorbox}[breakable,colback=teal!4,colframe=teal!35!black,title={Before you move on}]
What does \ja{まん} mean? (\textquotedblleft{}Ten thousand\textquotedblright{}.) Say it once more.
\end{tcolorbox}
\section[juuni]{\ja{じゅうに} (juuni) — twelve}

\label{lesson:JA-C60-juuni}



\begin{quote}
Before the new one: say the Japanese for a thousand, then the Japanese for ten thousand.
\end{quote}

\subsection*{You'll want to know: \ja{じゅうに}}
\textbf{\ja{じゅうに}} — \emph{juuni} — \textquotedblleft{}twelve\textquotedblright{}.

\begin{grammarlens}[title={What you've built}]
One more word for this chapter.
\end{grammarlens}

\subsection*{Your turn}
\begin{itemize}
\raggedright
  \item \emph{Say it:} \emph{juuni}
  \item \emph{Say it:} \emph{juuni}, once more
  \item \emph{Say it:} say \emph{juuni}
  \item \emph{Recall:} say the Japanese for a thousand, then the Japanese for ten thousand, then say \emph{juuni} again
\end{itemize}

\begin{tcolorbox}[breakable,colback=teal!4,colframe=teal!35!black,title={Before you move on}]
What does \ja{じゅうに} mean? (\textquotedblleft{}Twelve\textquotedblright{}.) Say it once more.
\end{tcolorbox}
\section[nasu]{\ja{なす} (nasu) — an aubergine}

\label{lesson:JA-C60-nasu}



\begin{quote}
Before the new one: say the Japanese for ten thousand, then the Japanese for twelve.
\end{quote}

\subsection*{You'll want to know: \ja{なす}}
\textbf{\ja{なす}} — \emph{nasu} — \textquotedblleft{}an aubergine\textquotedblright{}.

\begin{grammarlens}[title={What you've built}]
One more word for this chapter.
\end{grammarlens}

\subsection*{Your turn}
\begin{itemize}
\raggedright
  \item \emph{Say it:} \emph{nasu}
  \item \emph{Say it:} \emph{nasu}, once more
  \item \emph{Say it:} say \emph{nasu}
  \item \emph{Recall:} say the Japanese for ten thousand, then the Japanese for twelve, then say \emph{nasu} again
\end{itemize}

\begin{tcolorbox}[breakable,colback=teal!4,colframe=teal!35!black,title={Before you move on}]
What does \ja{なす} mean? (\textquotedblleft{}An aubergine\textquotedblright{}.) Say it once more.
\end{tcolorbox}
\section[uta]{\ja{うた} (uta) — a song}

\label{lesson:JA-C60-uta}



\begin{quote}
Before the new one: say the Japanese for twelve, then the Japanese for an aubergine.
\end{quote}

\subsection*{You'll want to know: \ja{うた}}
\textbf{\ja{うた}} — \emph{uta} — \textquotedblleft{}a song\textquotedblright{}.

\begin{grammarlens}[title={What you've built}]
One more word for this chapter.
\end{grammarlens}

\subsection*{Your turn}
\begin{itemize}
\raggedright
  \item \emph{Say it:} \emph{uta}
  \item \emph{Say it:} \emph{uta}, once more
  \item \emph{Say it:} say \emph{uta}
  \item \emph{Recall:} say the Japanese for twelve, then the Japanese for an aubergine, then say \emph{uta} again
\end{itemize}

\begin{tcolorbox}[breakable,colback=teal!4,colframe=teal!35!black,title={Before you move on}]
What does \ja{うた} mean? (\textquotedblleft{}A song\textquotedblright{}.) Say it once more.
\end{tcolorbox}
\section[eda]{\ja{えだ} (eda) — a branch}

\label{lesson:JA-C60-eda}



\begin{quote}
Before the new one: say the Japanese for an aubergine, then the Japanese for a song.
\end{quote}

\subsection*{You'll want to know: \ja{えだ}}
\textbf{\ja{えだ}} — \emph{eda} — \textquotedblleft{}a branch\textquotedblright{}.

\begin{grammarlens}[title={What you've built}]
One more word for this chapter.
\end{grammarlens}

\subsection*{Your turn}
\begin{itemize}
\raggedright
  \item \emph{Say it:} \emph{eda}
  \item \emph{Say it:} \emph{eda}, once more
  \item \emph{Say it:} say \emph{eda}
  \item \emph{Recall:} say the Japanese for an aubergine, then the Japanese for a song, then say \emph{eda} again
\end{itemize}

\begin{tcolorbox}[breakable,colback=teal!4,colframe=teal!35!black,title={Before you move on}]
What does \ja{えだ} mean? (\textquotedblleft{}A branch\textquotedblright{}.) Say it once more.
\end{tcolorbox}
